\section{Constraint-compatible numerical approximation}
\label{sec:numerical}

We consider a two-component instance of~\eqref{eq:main_system} on $Q_T=(0,0.25)\times(0,1)$:
\begin{align}
 u_t-d_1u_{xx}&=-u|u_x|^2, \label{eq:test-u}\\
 v_t-d_2v_{xx}&=u|u_x|^2-v|v_x|^2, \label{eq:test-v}
\end{align}
with $d_1=0.1$, $d_2=0.5$, homogeneous Dirichlet conditions, and
\[
 u(0,x)=0.8\sin^2(\pi x),\qquad v(0,x)=0.55\sin^2(2\pi x).
\]
This system has a transparent transfer--dissipation interpretation. The term $u|u_x|^2$ removes the first state and feeds the second at the same local rate, while $v|v_x|^2$ dissipates the second state. Thus steep gradients activate the chain $u\to v\to$ loss. This mechanism is relevant as a generic benchmark for front-sensitive conversion near concentration fronts, material interfaces, biological boundaries, or image edges. It is not asserted to be a calibrated model of any one of these phenomena; such a claim would require dimensional parameters and observational validation.

For the first component, the full preceding-gradient bound required by Lemma~\ref{lem:parametric-young} is not an additional assumption. Multiplication of~\eqref{eq:test-u} by $u$, followed by integration by parts, gives
\begin{equation}
 \frac12\|u(t)\|_{L^2(0,1)}^2
 +d_1\int_0^t\!\|u_x(s)\|_{L^2(0,1)}^2\,ds
 +\int_0^t\!\int_0^1u^2|u_x|^2\,dx\,ds
 =\frac12\|u_0\|_{L^2(0,1)}^2.
 \label{eq:u-energy}
\end{equation}
Consequently,
\[
 \int_0^T\|u_x(s)\|_{L^2(0,1)}^2\,ds
 \le \frac{\|u_0\|_{L^2(0,1)}^2}{2d_1}.
\]
Thus, for this two-component example, the hypothesis needed to control the $r=2$ cross term is verified directly. Equation~\eqref{eq:u-energy} also supplies a common diagnostic for the three numerical approximations.
This test is directly connected to the analysis: for $u,v\ge0$,
\[
 f_1=-u|u_x|^2\le0,\qquad f_1+f_2=-v|v_x|^2\le0,
\]
and $f_1(t,x,0,v,0,v_x)=0$, $f_2(t,x,u,0,u_x,0)=u|u_x|^2\ge0$. Thus the nonlinearities are quasi-positive, have triangular mass control, and exhibit quadratic gradient growth. Integrating the sum formally gives
\[
 \frac{d}{dt}\int_0^1(u+v)\,dx
 =d_1u_x\big|_0^1+d_2v_x\big|_0^1-\int_0^1v|v_x|^2\,dx,
\]
so homogeneous Dirichlet data provide a mass bound rather than an exact conservation law.

\subsection{Scope of the two-component benchmark}

The analytical estimate is formulated for arbitrary $r\ge2$, whereas the present benchmark has $m=r=2$. It verifies the common approximation protocol, the exact $r=2$ identity, the availability of the preceding-gradient bound, and the associated energy diagnostic. It does not numerically exercise the genuinely multi-cross-term situation $r\ge3$, in which the quantities $a_j=|d_r-d_j|$ must be summed simultaneously and the parametric choice in Lemma~\ref{lem:parametric-young} is most informative. The non-monotone three-coefficient calculation given in Section~\ref{sec:theoretical} illustrates that algebraic situation. Section~\ref{subsec:three-component} below now complements it with a three-component numerical benchmark, so that the simultaneous multi-cross-term case is no longer left purely at the algebraic level.
\subsection{Common comparison protocol}

All methods solve~\eqref{eq:test-u}--\eqref{eq:test-v} with identical coefficients and data. Results are evaluated at $101\times51$ common space--time points. A fine semi-implicit finite-difference computation with $401$ spatial points and $\Delta t=2.5\times10^{-5}$ is used as the numerical reference. Against an additional $801$-point computation with $\Delta t=1.25\times10^{-5}$, this reference has RMSE $8.62\times10^{-6}$ and relative $L^2$ error $2.96\times10^{-5}$.

The reported finite-difference method (FDM) uses $101$ spatial points and $\Delta t=2\times10^{-4}$. The Galerkin finite-element method (FEM) uses $100$ continuous piecewise-linear elements, the consistent mass matrix, two-point Gauss quadrature for the nonlinear element loads, and the same time step. Diffusion is treated implicitly and the nonlinear reaction explicitly in both methods.


\subsection{Manufactured-solution verification}
\label{subsec:mms}

Before comparing the three approximation paradigms, we verify the two classical discretizations independently of the refined-reference solution used later. This is a code-verification step in the sense of the Method of Manufactured Solutions (MMS)~\cite{roache2002mms,oberkampf2010vv}, rather than a validation of the physical prototype. The classical spatial and time discretizations are consistent with standard parabolic Galerkin and reaction--diffusion time-integration frameworks~\cite{thomee2006galerkin,hundsdorfer2003adr}. On $Q_{0.02}=(0,0.02)\times(0,1)$, prescribe
\[
 u_{\rm ex}(t,x)=0.8e^{-1.2t}\sin^2(\pi x),\qquad
 v_{\rm ex}(t,x)=0.55e^{-0.9t}\sin^2(2\pi x),
\]
which satisfy the homogeneous Dirichlet conditions. We add forcing terms
\begin{align*}
 s_1&=(u_{\rm ex})_t-d_1(u_{\rm ex})_{xx}+u_{\rm ex}|(u_{\rm ex})_x|^2,\\
 s_2&=(v_{\rm ex})_t-d_2(v_{\rm ex})_{xx}-u_{\rm ex}|(u_{\rm ex})_x|^2+v_{\rm ex}|(v_{\rm ex})_x|^2
\end{align*}
to~\eqref{eq:test-u}--\eqref{eq:test-v}. The derivatives used by the verification script are
\[
 (u_{\rm ex})_t=-1.2u_{\rm ex},\quad
 (u_{\rm ex})_x=0.8\pi e^{-1.2t}\sin(2\pi x),\quad
 (u_{\rm ex})_{xx}=1.6\pi^2e^{-1.2t}\cos(2\pi x),
\]
\[
 (v_{\rm ex})_t=-0.9v_{\rm ex},\quad
 (v_{\rm ex})_x=1.1\pi e^{-0.9t}\sin(4\pi x),\quad
 (v_{\rm ex})_{xx}=4.4\pi^2e^{-0.9t}\cos(4\pi x).
\]
The forcing expressions are generated from the displayed analytical derivatives. The evidential code-verification step is the independent refinement study below, rather than a zero obtained by re-evaluating the same defining residual. Spatial and temporal refinements are performed separately. The spatial studies use $\Delta t=2\times10^{-6}$ for FDM and FEM. The temporal study uses the steps $8,4,2,1\times10^{-4}$ with $801$ FDM points and $300$ $P_1$ finite elements, so that the first-order IMEX time error dominates the remaining spatial error. This test verifies the implementations; it is not a claim about convergence of the unforced nonlinear problem under arbitrary data.

\subsection{Three-component non-monotone benchmark}
\label{subsec:three-component}

To exercise the simultaneous cross-term configuration, we additionally solve
\begin{align}
 u_t-d_1u_{xx}&=-u|u_x|^2,\label{eq:three-u}\\
 v_t-d_2v_{xx}&=-v|v_x|^2,\label{eq:three-v}\\
 w_t-d_3w_{xx}&=u|u_x|^2+v|v_x|^2-w|w_x|^2,\label{eq:three-w}
\end{align}
with
\[
 (d_1,d_2,d_3)=(0.1,0.8,0.3),
\]
zero Dirichlet data, and
\[
 u_0=0.8\sin^2(\pi x),\qquad
 v_0=0.55\sin^2(2\pi x),\qquad
 w_0=0.35\sin^2(3\pi x).
\]
For nonnegative states,
\[
 f_1\le0,\qquad f_1+f_2=-u|u_x|^2-v|v_x|^2\le0,
 \qquad f_1+f_2+f_3=-w|w_x|^2\le0.
\]
Quasi-positivity is explicit: $f_1(0,\cdot)=0$, $f_2(0,\cdot)=0$, and $f_3(u,v,0)=u|u_x|^2+v|v_x|^2\ge0$. The nonlinearities also have the quadratic gradient growth required by the structural setting of Section~\ref{sec:theoretical}; for example $|u||u_x|^2\le C(|u|)(1+|u_x|^2)$, and the remaining terms satisfy analogous bounds. The benchmark is deliberately designed as a structural verification problem rather than as a calibrated application: it exposes two simultaneous cross terms while keeping the preceding-gradient estimates independently auditable. Moreover, multiplying~\eqref{eq:three-u} and~\eqref{eq:three-v} by their own components gives
\[
 C_1\le \frac{\|u_0\|_2^2}{2d_1}
 =\frac{0.8^2(3/8)}{0.2}=1.2,
\qquad
 C_2\le \frac{\|v_0\|_2^2}{2d_2}
 =\frac{0.55^2(3/8)}{1.6}=0.0708984375.
\]
Thus Assumption~\ref{ass:previous-full-gradient} is independently verified for both preceding components. The same bounds hold uniformly for the sign-preserving cut-off approximants of Section~\ref{sec:theoretical}: multiplication by the common nonnegative cut-off changes the dissipative terms to $-\chi_n u^2|u_x|^2$ and $-\chi_n v^2|v_x|^2$, which retain their sign. Hence the estimates needed in Proposition~\ref{prop:partial-sum} are available at the approximation level, not only for the formal benchmark equations. Since $a_1=0.2$, $a_2=0.5$, $A_3=0.7$ and $\varepsilon=d_3/A_3=3/7$, Lemma~\ref{lem:parametric-young} gives
\[
 \frac{a_1A_3}{2d_3}C_1+\frac{a_2A_3}{2d_3}C_2
 =\frac7{30}C_1+\frac7{12}C_2
 \le 0.321357421875.
\]
The numerical reference uses $401$ FDM points and $\Delta t=2.5\times10^{-5}$. A separate $801$-point computation with $\Delta t=1.25\times10^{-5}$ differs from it by RMSE $1.03\times10^{-5}$ and relative $L^2$ error $4.24\times10^{-5}$. Reported FDM and FEM runs use, respectively, $101$ spatial points and $100$ linear elements with $\Delta t=2\times10^{-4}$. The three-component experiment is used as a structural benchmark for the analytical estimate. The PINN comparison is intentionally confined to the two-component problem so that this coercivity-focused experiment is not conflated with an additional neural-optimization study.

\subsection{PINN formulation}

The PINN has four hidden layers of $48$ tanh neurons. Its output is parameterized as
\[
 (u_\theta,v_\theta)(t,x)=x(1-x)\operatorname{softplus}(N_\theta(t,x)),
\]
which enforces homogeneous boundary conditions and non-negativity exactly. Exact output transformations are attractive because penalty-based Dirichlet enforcement can require delicate loss-weight tuning~\cite{berrone2023dirichlet}. The loss is
\[
 \mathcal L(\theta)=\frac1{N_f}\sum_{q=1}^{N_f}
 \left(|R_{1,\theta}(t_q,x_q)|^2+|R_{2,\theta}(t_q,x_q)|^2\right)
 +20\frac1{N_0}\sum_{q=1}^{N_0}|U_\theta(0,x_q)-U_0(x_q)|^2,
\]
where $R_{1,\theta}$ and $R_{2,\theta}$ are the residuals of~\eqref{eq:test-u}--\eqref{eq:test-v}. The initial-condition weight $20$ is fixed a priori for all reported runs and is not tuned against the refined reference; no claim of optimality is attached to this choice. We use $N_f=3000$, $N_0=200$, Adam with learning rate $10^{-3}$, and $2000$ epochs. Five independent runs use seeds 11, 29, 47, 71, and 101. A common independently sampled set of 5000 points, generated with seed 2025, is used to evaluate every trained network. The realization whose RMSE is closest to the five-seed median (seed 29) is used in field figures; all aggregate PINN values use the five runs. Automatic differentiation evaluates the derivatives of the neural representation; it does not remove approximation, optimization, or floating-point error.
The initial condition is imposed by a penalty deliberately. The commonly suggested map $U_\theta=U_0+t\,x(1-x)\operatorname{softplus}(N_\theta)$ would enforce initial and boundary data together with non-negativity, but it would also impose $U_\theta(t,x)\ge U_0(x)$ pointwise and is therefore incompatible with the dissipative decay observed here. The signed alternative $U_0+t\,x(1-x)N_\theta$ avoids that monotonicity constraint but no longer guarantees non-negativity. We consequently retain the nonnegative boundary-compatible output map and quantify the remaining mismatch explicitly through the relative initial-data error $E_{\mathrm{IC}}$ and the mass diagnostic. Constructing an ansatz that enforces all three constraints without restricting the temporal dynamics is a separate approximation-design problem.

\subsection{Metrics and reproducibility}

For a method $A$, we report
\[
 \operatorname{RMSE}_A=\left(\frac1{2N}\sum_{q=1}^{N}
 \bigl[(u_A-u_{\rm ref})^2+(v_A-v_{\rm ref})^2\bigr]\right)^{1/2}
\]
and the relative discrete $L^2$ error
\[
 E_{\mathrm{rel},A}=\frac{\left(\sum_{q=1}^{N}[(u_A-u_{\rm ref})^2+(v_A-v_{\rm ref})^2]\right)^{1/2}}{\left(\sum_{q=1}^{N}[u_{\rm ref}^2+v_{\rm ref}^2]\right)^{1/2}}.
\]
Constraint diagnostics are the minimum value over interior evaluation points and the maximum positive increment of the trapezoidal approximation to $\int_0^1(u+v)\,dx$. The PINN initial-data mismatch is additionally reported as
\[
E_{\mathrm{IC}}=\frac{\left(\sum_q[(u_\theta(0,x_q)-u_0(x_q))^2+(v_\theta(0,x_q)-v_0(x_q))^2]\right)^{1/2}}{\left(\sum_q[u_0(x_q)^2+v_0(x_q)^2]\right)^{1/2}}.
\] Training or solve time and PINN inference time are implementation-dependent diagnostics and are not interpreted as hardware-normalized complexity measurements. Python-traced memory is retained only in the raw files because it excludes native TensorFlow/SciPy allocations and is not used for scientific comparison. The PINN differential residual is also evaluated on the common set of $5000$ independently sampled points that are never used for training. Training and test residuals are both reported as root-mean-square quantities. PINN statistics use the sample standard deviation; with five runs, they are descriptive variability measures rather than confidence intervals.

The supplementary reproducibility archive contains all source files, parameters, seeds, raw metrics, loss histories, metadata, and figure-generation routines. It provides both a quick pipeline check and the publication configuration. The archived run used Python 3.12.10, NumPy 2.5.2, SciPy 1.18.0, and TensorFlow 2.21.0 on Windows 11 with an Intel64 processor. FDM/FEM use float64 and the PINN uses float32; timings remain implementation-specific. The archive distributed with the manuscript is the reproducibility source of record for the reported calculations.